\documentclass[11pt,a4paper]{article}
\usepackage[margin=22mm]{geometry}
\usepackage{amsmath,amssymb,booktabs,microtype}
\usepackage[hidelinks]{hyperref}
\newcommand{\ii}{\mathrm i}
\newcommand{\tr}{\operatorname{tr}}
\newcommand{\ket}[1]{|#1\rangle}
\newcommand{\bra}[1]{\langle#1|}
\title{Route G-R9: One Relational Spectrum, Two Dials,\\and a Missing Exchange Channel}
\author{A finite candidate and its failure boundary; no mass or time-law fit}
\date{9 October 2026}
\begin{document}
\maketitle
\begin{abstract}
Test the event/spectrum interpretation before importing laboratory
spacetime. A three-node connection graph supplies one dimensionless
spectrum. Two independent copies supply two transition dials. A fixed
selected-mode interaction gives exact resonant exchange, but a fixed
vertex contact opens an additional degenerate channel and invalidates
the ideal two-channel closure. This is a useful rejection of a hidden
selection assumption, not a universal no-go result. Ordered updates,
quantum mechanics, preparation and readout remain inputs: neither
physical time nor particle masses are derived.
\end{abstract}

\section{Starting point: relations before coordinates}
Chen's sections 1.2, 1.8 and 1.9 motivate event accumulation, a visible
spectrum and exchange between compatible event rhythms~\cite{chen}.
They do not specify the matrices below. We supply complex amplitudes,
Born probabilities and a three-node graph; nodes are not three physical
dimensions. The phase below is a link phase, not a rotation of matter.

Let $S\ket j=\ket{j+1\pmod3}$ and define
\begin{align}
 B_\theta&=I-e^{\ii\theta/3}S,&
 L_\theta&=B_\theta^\dagger B_\theta
 =2I-e^{\ii\theta/3}S-e^{-\ii\theta/3}S^\dagger,\label{L}\\
 Q_\theta(z)&=\|B_\theta z\|^2
 =\sum_{j=0}^2|z_{j+1}-e^{\ii\theta/3}z_j|^2\geq0.
\end{align}
The normalized vectors $u_k(j)=3^{-1/2}e^{2\pi\ii kj/3}$ obey
$Su_k=e^{-2\pi\ii k/3}u_k$. Thus, without a fitted eigenvalue,
\begin{equation}
 L_\theta u_k=\lambda_k(\theta)u_k,\qquad
 \lambda_k=2-2\cos\left(\frac{\theta-2\pi k}{3}\right).\label{spec}
\end{equation}
For $0<\theta<\pi$, $\lambda_0<\lambda_1<\lambda_2$; endpoints can
be degenerate. In particular,
\begin{equation}
 \operatorname{spec}L_0=\{0,3,3\},\qquad
 \operatorname{spec}L_{\pi/2}=\{2-\sqrt3,2,2+\sqrt3\}.
\end{equation}
All diagonal entries remain 2, all edge magnitudes remain 1 and
$\tr L_\theta=6$. Scalar averages do not determine the full spectrum.
We do \emph{not} identify $\tr L/3$ with the source's undefined $E_d$.

\newpage
\section{What is invariant, stable, and still assumed}
A local basis rephasing $D=\operatorname{diag}(e^{\ii\chi_j})$ sends
$L\mapsto DLD^\dagger$ without changing eigenvalues. With oriented
hopping amplitudes $a_{j+1,j}=-L_{j+1,j}$,
\begin{equation}
 a_{10}a_{21}a_{02}=e^{\ii\theta}
\end{equation}
is unchanged by $D$. Changing this loop product is not a coordinate
change. Reversing loop orientation changes the sign convention for
$\theta$, not the unordered spectrum.

The normalized lowest mode minimizes $Q$ on $\|z\|=1$ and is isolated
by $d=\lambda_1-\lambda_0>0$. For Hermitian $E$ with
$\|E\|=\varepsilon<d/2$, the min--max principle gives
\begin{equation}
 |\lambda'_j-\lambda_j|\leq\varepsilon,\qquad
 \lambda'_1-\lambda'_0\geq d-2\varepsilon.
\end{equation}
Writing $P_\perp=I-u_0u_0^\dagger$, project the perturbed eigen-equation:
$(P_\perp LP_\perp-\lambda'_0)P_\perp u'_0=-P_\perp Eu'_0$.
Its left operator has inverse norm at most $(d-\varepsilon)^{-1}$, so
\begin{equation}
 \|P_\perp u'_0\|\leq\frac{\varepsilon}{d-\varepsilon}.
\end{equation}
This is spectral stability, not dynamical attraction or a stable particle.

\paragraph{The extra dynamical postulate.}
Choose one update rule $U=e^{-\ii\alpha H}$, with $\alpha=1/4$ fixed
before varying $\theta$, and $\Psi_{n+1}=U\Psi_n$. The ordered integer
$n$ counts \emph{assumed} updates; it is neither measured seconds nor
an order derived from the graph. For a free copy $H=L$, every $u_k$
persists as a ray under these updates. A compatible finite history
minimizes $\sum_n\|\Psi_{n+1}-U\Psi_n\|^2$ at zero for a given input.
This positive history functional supplies no event-ordering principle.
We use $s=\alpha n$ as an interpolation variable only.

\paragraph{Two copies, not two qubits inside one qutrit.}
Take $\mathcal H_A\otimes\mathcal H_B=\mathbb C^3\otimes\mathbb C^3$.
Use levels $(u_0,u_1)$ for A and $(u_1,u_2)$ for B. The gap pair is
\begin{equation}
 \Delta_A=\lambda_1-\lambda_0,\qquad
 \Delta_B=\lambda_2-\lambda_1.\label{gaps}
\end{equation}
For each free plus-state dial, $c=2\rho_{10}=e^{-\ii\Delta s}$.
Two fixed quadratures determine its real and imaginary parts in
separate preparations. With conventional Pauli $Y$, $\langle X\rangle
=\Re c$ and $\langle Y\rangle=\Im c$. Where phases can be unwrapped,
$\phi_B=(\Delta_B/\Delta_A)\phi_A$ eliminates $s$. This alone does
not produce saved records, an epoch counter or an autonomous clock.

\newpage
\section{An exactly solvable, but selected, exchange rule}
In the fixed Fourier basis (independent of $\theta$), postulate
\begin{equation}
 H=L\otimes I+I\otimes L+
 g\big(\ket{1,1}\bra{0,2}+\ket{0,2}\bra{1,1}\big),\quad g=0.1.
 \label{selected}
\end{equation}
This selection is a physical assumption, not a consequence of $L$ or
nearest-edge locality. For this rule only, the four-dimensional sector
$(\ket{0,1},\ket{0,2},\ket{1,1},\ket{1,2})$ is exactly invariant.
Subtract the common scalar $\lambda_0+\lambda_1$ and use logical labels
$(\ket{00},\ket{01},\ket{10},\ket{11})$:
\begin{equation}
 H_4=\begin{pmatrix}
 0&0&0&0\\0&\Delta_B&g&0\\0&g&\Delta_A&0\\
 0&0&0&\Delta_A+\Delta_B
 \end{pmatrix}.\label{H4}
\end{equation}
No rotating-wave truncation was made in this particular model. Its
one-excitation block has average $\bar\Delta=(\Delta_A+\Delta_B)/2$
and splitting $\Omega=\sqrt{(\delta/2)^2+g^2}$, where
$\delta=\Delta_A-\Delta_B$. Exponentiating its traceless square
$[(H_4-\bar\Delta I)_{1\rm exc}]^2=\Omega^2 I$ yields
\begin{equation}
 P_{10\to01}(n)=\frac{g^2}{\Omega^2}\sin^2(\Omega\alpha n),\qquad
 P\leq\frac{4g^2}{\delta^2+4g^2}.\label{exchange}
\end{equation}
The upper envelope is over continuous interpolation; an integer
update need not attain it. Conservation here is of total $H$ and of
the selected total excitation number, not of the separate free clocks.

Using $\lambda_0+\lambda_1+\lambda_2=6$,
\begin{equation}
 \delta=3\lambda_1-6=-6\cos\left(\frac{\theta-2\pi}{3}\right).
\end{equation}
Thus exact one-for-one resonance in $0<\theta<\pi$ occurs only at
$\theta=\pi/2$; both gaps are $\sqrt3$. Near it
$\delta=-2(\theta-\pi/2)+O((\theta-\pi/2)^3)$.
Nothing makes this phase an attractor. Rational frequency ratios alone
do not imply exchange through this operator; higher-order common-multiple
processes require additional channels. Finite coupling also permits
off-resonant exchange, not a strict event-coincidence rule.

\begin{center}
\begin{tabular}{rrrrr}\toprule
$\theta/\pi$&$\Delta_A$&$\Delta_B$&free B/A&exchange envelope\\\midrule
$1/3$&2.226682&1.184793&0.532089&0.035539\\
$1/2$&1.732051&1.732051&1&1\\
$2/3$&1.184793&2.226682&1.879385&0.035539\\\bottomrule
\end{tabular}
\end{center}
At the resonant integer update $n=63$, $P=0.99998233$, not exactly 1.
These are model-unit predictions, not observed probabilities or a fit.

\newpage
\section{Clock visibility and the underdetermined update law}
The free gap ratio is not automatically an interacting clock rate.
For a product plus-state input at resonance, Eq.~\eqref{H4} gives
\begin{equation}
 \Psi(s)=\frac12\left[\ket{00}
 +e^{-\ii(\Delta+g)s}(\ket{01}+\ket{10})
 +e^{-2\ii\Delta s}\ket{11}\right],\quad \Delta=\sqrt3.
\end{equation}
Tracing either other copy gives
\begin{equation}
 c_A(s)=c_B(s)=e^{-\ii\Delta s}\cos(gs).\label{dials}
\end{equation}
The phases agree when readable, but visibility $|\cos(gs)|$ changes.
At its zeros neither reduced dial has a phase; crossings can give
$\pi$ phase jumps. A negative cosine is not reversed physical time.
The continuous zeros need not be sampled at integer $n$. Away from
resonance, exponentiating Eq.~\eqref{H4} and taking the partial traces
is required: the free formula is no longer exact. No unique rigid
relative clock rate is claimed in that case.

Exchange probabilities in Eq.~\eqref{exchange} use the different input
$\ket{10}$; dial coherences use $\ket{+,+}$. These are separate runs of
the same generator, not simultaneously sharp phase and excitation
readings. Unitary recurrences are not dissipative synchronization.
Neither protocol includes an autonomous preparation/measurement memory.

\paragraph{A scalar average is insufficient; $E_d$ remains undefined.}
Across the table, A slows per update while B speeds up, despite constant
edge magnitudes and trace. This rejects a sole common fractional factor
for these two transitions in this candidate. It does not refute an
$E_d$ containing relational phases or additional fields, and it does
not prove anything about an external clock's seconds.

\paragraph{The graph does not select the dynamical generator.}
Both $L$ and $f(L)=L+\eta L^2$, $\eta>0$, are positive with identical
eigenvectors. Keeping $\eta=0.2$ fixed, at the \emph{same} $\theta=\pi/2$,
\begin{equation}
 \frac{f(\lambda_2)-f(\lambda_1)}{f(\lambda_1)-f(\lambda_0)}
 =\frac{1+\eta(4+\sqrt3)}{1+\eta(4-\sqrt3)}=1.4766271\ldots,
\end{equation}
instead of 1. This is a different dynamical postulate, not a harmless
change of units. Stability and positivity alone cannot decide between
the two. A first-principles completion must supply an independent
principle for the rule, rather than choose it to obtain desired clocks.

\paragraph{Visibility is not mass conversion.}
A fixed normalized probe $q$ has the spectral measure
\[
 \sum_k|q^\dagger u_k|^2\delta(x-\lambda_k).
\]
A vertex probe gives weights $1/3$; the $u_0$ probe sees only its mode.
Changing the probe changes these weights, not the eigenvalues.
No Lorentzian dispersion, calibrated energy or propagator mass pole
has been supplied, so neither these eigenvalues nor their gaps are
identified with elementary-particle masses.

\newpage
\section{Fixed vertex contact: the missing-channel countercheck}
The selected rule in Eq.~\eqref{selected} is not the natural contact
\begin{equation}
 V_{\rm loc}=\gamma P_0\otimes P_0,\qquad
 P_0=\ket{j=0}\bra{j=0},\quad\gamma>0.\label{contact}
\end{equation}
``Local'' here means one node of each graph, not locality in derived
spacetime. Since every $u_k(0)=1/\sqrt3$, every matrix element of
$V_{\rm loc}$ in the two-copy Fourier basis is $v=\gamma/9$.
Let $E$ embed the selected four states and $P=EE^\dagger$. Then
\begin{equation}
 \|(I-P)V_{\rm loc}E\|=\sqrt{20}\,\gamma/9>0.
\end{equation}
Consequently that sector is not invariant. At $\theta=\pi/2$ the
three states $\ket{0,2},\ket{1,1},\ket{2,0}$ all have free total
eigenvalue 4. The third is not a high-energy state that can be dropped.

There is an exact bound independent of perturbation theory. The full
nine-dimensional generator commutes with copy swap. The input
$\ket{1,1}$ is symmetric, hence at every update
\begin{equation}
 P_{11\to02}=P_{11\to20},\qquad P_{11\to02}\leq\frac12.\label{half}
\end{equation}
This contradicts nearly unit single-channel transfer for this contact,
without tiny-error arguments. It does not exclude other interactions.

In the weak-contact resonant block, $H_{\rm res}=4I+v\mathbf1\mathbf1^T$,
where $\mathbf1=(1,1,1)^T$. Its symmetric vector has shift $3v$ and the
orthogonal two-dimensional subspace has shift zero. Therefore
\begin{equation}
 e^{-\ii s v\mathbf1\mathbf1^T}
 =I+\frac{e^{-3\ii vs}-1}{3}\mathbf1\mathbf1^T,
 \qquad P_{11\to02}^{\rm res}=P_{11\to20}^{\rm res}
 =\frac49\sin^2(3vs/2).\label{49}
\end{equation}
The single-target envelope is $4/9$; both targets together reach $8/9$.
Unlike Eq.~\eqref{half}, $4/9$ is not asserted as an exact bound on the
full generator. Other free multiplets are at distance $d=\sqrt3$.
Finite-dimensional degenerate perturbation gives block corrections
$O(\gamma^2/d)$ and spectral-projector changes $O(\gamma/d)$.
For bounded $\gamma s$, probabilities approach Eq.~\eqref{49} with
errors $O(\gamma/d)$. No long-time uniform claim is made.

\begin{center}
\begin{tabular}{rrrrr}\toprule
$\gamma$&update $n$&exact $P_{02}$&exact $P_{20}$&outside triple\\\midrule
0.09&419&0.444307380&0.444307380&$2.055\,10^{-4}$\\
0.03&1257&0.444415403&0.444415403&$4.346\,10^{-5}$\\
0.01&3770&0.444442499&0.444442499&$2.918\,10^{-6}$\\\bottomrule
\end{tabular}
\end{center}
These full-matrix fixtures use $n$ nearest the first secular peak;
off-resonant leakage is retained, not fitted away.

\newpage
\section{Decision, verification, and a bounded next step}
\paragraph{What this stage closes.}
A single positive graph rule supplies an explicit spectral structure,
phase invariants and persistent modes. Fixed transition choices on
two independent copies supply comparable dials. One selected interaction
has an exact resonance formula, while a vertex contact supplies a
counterexample to treating that selection as automatic. The latter is
the material new constraint: an omitted \emph{degenerate} channel, not
insufficient decimal accuracy, changes the interpretation.

\paragraph{What remains open.}
This is not a coordinate-free derivation of time: the update order,
quantum kinematics, preparation, reference quadratures and interaction
are assumed. No physical seconds, density-time equation, three-dimensional
space, Lorentzian propagation or particle mass is generated. The modes
are persistent rays, not self-organized attractors. G-R1's stationary
constraint and G-R2's recording apparatus are not silently appended to
this new generator. Their compatibility would require a separate test.

\paragraph{Research choice.}
Retain the full three-channel local-contact model as the economical
candidate. Next derive its two-copy observation algebra and conditional
readout, retaining unobserved channels, and test whether an operational
relative clock still exists. Alternatively, justify a symmetry or a
microscopic source that forbids the unwanted channel \emph{before}
using the selected-mode model. Do not rescue the ideal transfer by
deleting its degenerate partner or by selecting particle-specific rules.
Neither path needs more gravitational-interferometer calibration first.

G-R4--G-R8 remain useful thermal/metric controls, not derivations of
the original event concept and not blockers of this mainline. No
Route-F fitting, determinant, or portal gate is promoted.

\paragraph{Executable verification.}
The companion program is
\texttt{verify\_gr9\_event\_network.py}.
Its inputs are fixed in \texttt{EVENT\_NETWORK\_CARD.json}.
The 342 checks cover the Fourier
diagonalization, positive factorization, gauge and permutation invariance,
gap and eigenvector bounds, exact site/sector equivalence, repeated
integer updates, exchange and readout identities, the alternative
generator and the contact countercheck. The gap-bound checks use
$\pi/6\leq\theta\leq5\pi/6$ away from degeneracy. Deterministic fixtures are
not data fitting. Full $9\times9$ contact evolution is compared with,
not replaced by, the weak-coupling block. The report records all inputs
and preserves selected G-R1, G-R2 and G-R8 hashes. Numerical checks
support the finite algebra; they are not experiments or proofs of
the stated emergence claims.

\begin{thebibliography}{9}
\bibitem{chen} BoYu Chen, \emph{Energy spectrum for elementary particles},
final version, 23 April 2022, sections 1.2, 1.8 and 1.9.
\href{https://indico.cern.ch/event/1109513/contributions/4969967/attachments/2480973/4410560/Energy\%20spectrum\%20for\%20elementary\%20particles_final.pdf}{Author's five-page contribution on CERN Indico}.
Conceptual source, not evidence for the matrix dynamics in this note.
The locally audited PDF hash is recorded in the input card. The public
web fetch returned 403 in this run; the previously acquired, hash-pinned
copy was used. These elementary spectral/exchange identities are not
claimed as new mathematical theorems.
\end{thebibliography}
\end{document}
